% critique.tex — f163r, vv. 43–79 : commentaire critique
% Articles tirés du commentaire antérieur de l'auteur (texte et notes conservés) ; répartition : repartition.md
\cc{49–50}{μέλλει γὰρ ὁ καλλιεπὴς Ἀγάθων / πράμος ἡμέτερος.}{πράμος n'est ni autrement attesté ni même relevé par les grammairiens et doit nécessairement être scandé $\cup \cup $ pour rentrer dans le mètre anapestique, empêchant tout rapprochement étymologique pertinent. 
Autant d'indices qui font soupçonner une erreur textuelle plutôt qu'un hapax agathonien.
πρόμος, \italique{le chef}, fut proposé par Grynaeus et quasi-unanimement repris depuis. 
La correction est très légère mais n'est pas satisfaisante : 
il manque au terme πρόμος une acception \italique{sensu obsceno} pour expliquer l'exclamation de Mnésiloque : μῶν βινεῖσθαι;\nc{Sommerstein voit en πρόμος une légère déformation suggérant à l'oreille πόρνος... « Notre chef de file --- Celui qu'on enfile ? » Mais que c'est loin en grec !.. Certains ont voulu justifier le rapprochement πρόμος -- πόρνος en arguant de la plaisanterie de \italique{Nuées} (v. 394) sur l'indissociabilité entre βροντή et πορδή -- certes pas au niveau de la prononciation. J'ai profité des annexes auxquelles je renvoie le lecteur pour exégéter ce jeu de mots qui roule sur la langue proverbiale.} 
De plus, l'on s'attendrait plutôt à trouver un infinitif réclamé par μέλλει. Une glose d'Hésychius va nous permettre de débrouiller l'écheveau. Cette glose est elle-même légèrement corrompue (π 3188) :\citComm{†πραῖνοι†· πρηνίζειν, καταστρέφειν.}Sans nous attarder pour l'instant sur les hypothèses de corruptions, constatons que Hésychius atteste l'existence du verbe πρᾳνόω (ou πρανόω\nc{Même genre de variantes orthographiques entre πρᾷος et πρᾶος.}).
Pour le sens, πρηνίζειν signifie \italique{précipiter}, \italique{renverser}, \italique{jeter} (en avant).
Le serviteur d'Agathon construit ainsi sa phrase :\citComm{μέλλει γὰρ Ἁγάθων πρανοῦν δρυόχους\nc{Les δρύοχοι sont les \italique{chantiers}, ou poutres sur lesquelles s'établit la quille d'un vaisseau en construction. On rapprochera la métaphore de \italique{Guêpes} 30 : τὴν τρόπιν τοῦ πράγματος, « la base de l'affaire ». (Note de Taillardat §758 n. 1 ; voir aussi §673.)}.\\
Agathon s'apprête à jeter les fondations.}Seulement, πρηνίζειν peut désigner aussi l'action de s'accroupir tête baissée\nc{Cf. Hésychius π 3199, π 3269 πρηνιχθέντα· ἐπὶ πρόσωπον καὶ ἐπὶ στόμα πεσόντα ~ Photius π 1138 et l'Etymologicum genuinum : πρηνίξαι πρηνιχθῆναι καὶ πεπρηνίχθαι· γνὺξ γάρ ἐστιν κυρίως τὸ εἰς γόνυ πεσεῖν ... τὸ δὲ πρηνιχθῆναι ἐπὶ τὸ πρόσωπον [πεσεῖν] et d'autres références dans l'abondant et récent apparat des témoins de l'édition Theodoridis de Photius.}. (πρανοῦν, comme accroupir, fait entendre la \italique{croupe}, πρανής.) C'est bien-sûr dans ce sens détourné que le prend Mnésiloque -- image d'un Agathon se jetant à son travail dans la position du penseur de Rodin ? L'occasion est trop belle : Mnésiloque s'imagine Agathon en position κύβδα (cf. le vers 489) et tournant le dos à d'éventuels assaillants --- μῶν βινεῖσθαι; demande-t-il. Le véritable Agathon devait effectivement travailler dans une posture qui prêtait à rire aux mauvais esprits : c'est sur son lit et dans cette même posture qu'il apparaît sur scène quelques vers plus loin ; cette position étrange lui vaudra d'être comparé par Mnésiloque à Cyrène \italique{aux douze postures} (δωδεκαμήχανος) (v. 98).\\
Reste à expliquer la présence dans le \italique{Ravennas} de †ἡμέτερος† après †πράμος†. Hésychius glose †πραῖνοι† par deux infinitifs, πρηνίζειν et καταστρέφειν, preuve que le mot relevé en vedette devait originellement être lui aussi un infinitif. Le rétablissement de πρανοῦν se fait alors sans difficulté. C'est la même mécoupure du verbe πρανοῦν, lu †πρανοι \ ν† qui pourrait expliquer le †πράμος† et le début de †ἡμέτερος† du \italique{Ravennas} selon le \italique{stemma} suivant :  \begin{align*}
\overset{\overset{\text{\footnotesize{\allemand{Urtext}}}}{\text{πρηνίζειν, καταστρέφειν}}}{\text{πρανοῦν [...]}}\xrightarrow{\text{mécoup. \ } + \text{\ conf. \  υ/ι}}\overset{\overset{\text{\footnotesize{Archétype}}}{\text{πρηνίζειν, καταστρέφειν}}}{\text{πρανοι \ \ ν[...]}}&\xrightarrow{\text{\ \ \ \ \ }}\overset{\text{\footnotesize{Hésychius π 3188}}}{\text{πραῖνοι· πρηνίζειν, καταστρέφειν}}\\
& \xrightarrow{\text{confus. \ ν/μ et ι/ς}}\overset{\text{\footnotesize{\italique{Ravennas}}}}{\text{πραμος <η>μ[ετερος]}}
\end{align*}
Rien de plus semblable qu'un ι et un ς, ainsi qu'un μ et un ν  -- dans les manuscrits écrits en minuscules à l'instar du \italique{Ravennas}.
Il s'agit maintenant de comprendre la présence de †ετερος†. Sans tenter d'explications, je constate qu'au vers 562 des \italique{Nuées}, les deux principaux manuscrits des pièces d'Aristophane, notre \italique{Ravennas} et le \italique{Venetus}, portent deux leçons différentes : ἑτέρας \italique{versus} προτέρας\nc{Les écrits anciens comportaient de nombreuses abréviations. (Cf. par exemple la très ingénieuse reconstruction du vers 80 par Nauck.) Par exemple, προσ était souvent abrégé d'une sorte de Ʒ renversé qui ressemblait à un epsilon : \renverse{Ʒ}. Voir G. Bady, \italique{Paléographie grecque : éléments de cours}, ou bien V. Gardthausen, \italique{Griechische palaeographie}\expo{2}, II (1913), p. 340.}. 
Je suppose ici la même corruption :\citComm{μέλλει γὰρ ὁ καλλιεπὴς Ἀγάθων πρανοῦν προτέρους δρυόχους, τιθέναι δράματος ἀρχάς. κάμπτει δὲ νέας ἁψῖδας ἐπῶν κτλ.\\
Agathon au beau langage s'apprête à jeter les premiers fondements, à disposer les prémices d'un drame. Il tourne de nouvelles roues de mots etc.}}
\cc{61}{γογγυλίσας καὶ συστρέψας.}{Vers retouché \italique{gratia metri} d'après un passage de \italique{Lysistrata} (v. 975), jouée probablement l'année précédente (cf. Introduction).}
\cc{70–4}{τί οὖν ἐγὼ δρῶ; ... ὄντα κηδεστὴν ἐμόν.}{La répartition de ce passage diffère d'un éditeur à un autre ; chacun voulant apporter son grain de sel. Je m'inscris, moi aussi, dans cette grande tradition philologique et j'apporte ma pierre à l'édifice : une nouvelle répartition des répliques qui, je l'espère, clôturera un débat qui n'a que trop duré. La répartition des deux répliques du vers 70 ne pose aucun problème : il suffit de songer aux caractères des trois personnages en présence : le ton empoté du « τί οὖν ἐγὼ δρῶ; » ne peut convenir qu'au personnage pour l'instant passif et qui subit l'action : Mnésiloque\nc{Réplique analogue de Mnésiloque au vers 925 : ἐγὼ δ᾽ ὁ κακοδαίμων τί δρῶ;} (Euripide ruisselle de défauts mais garde toujours une certaine dignité ; jamais il ne s'abaisserait à lancer un mot aussi trivial et décontenancé – ses plaintes ont davantage de lyrisme !) La réponse : « περίμεν᾽, ὡς ἐξέρχεται » ne peut être attribuée qu'à Euripide : c'est lui qui a besoin de Mnésiloque ; de plus, le serviteur ne délivre ni conseil ni ordre : il se contente de répondre le plus sèchement possible aux questions (cette répartition est d'ailleurs celle du manuscrit de Ravenne). Les vers 72 à 74 sont évidemment prononcés par Mnésiloque qui cherche à percer la détresse d'Euripide et nous révèle par la même occasion son lien de parenté avec le poète. Reste le vers 71 : il est plus que tentant de l'attribuer à Mnésiloque et faire ainsi des quatre vers 71-74 l'exact pendant des quatre vers 1-4 : une invocation à Zeus suivie d'une question : là pour savoir où Euripide le mène et ici pour savoir ce qui le tourmente. Problème : la présence d'un \italique{signe} (σημεῖον), probablement un \italique{paragraphos}, sur d'anciens manuscrits a été interprétée comme indiquant un changement d'interlocuteurs entre le vers 71 et les vers 72-74. Mais tel n'est pas le cas ici, et c'est ce que nous indiquait une ancienne scholie, très précieuse et incomprise\nc{Incomprise déjà de Suidas – lexicographe hypothétique du \nieme{IX} ou \nieme{X} siècle, à qui l'on attribue le lexique de la \italique{Souda}. Ce dernier fourmille en références et citations d'Aristophane : la \italique{Souda} est un outil de recherche sur Aristophane indispensable. 
Il en existe une version numérique : \href{https://www.cs.uky.edu/~raphael/sol/sol-html/}{https://www.cs.uky.edu/\textasciitilde raphael/sol/sol-html/} – incomprise donc de Suidas ce qui tend à prouver son ancienneté : Suidas (κ 1487) biffe le lemme `τὸ σημεῖον', pourtant indispensable à la bonne compréhension et égare ainsi les différents éditeurs des scholies qui l'appliquent au vers 74, voire qui la retouchent pour finir de la rendre tout-à-fait inintelligible.} : 
\citComm{τὸ σημεῖον· οὐ μόνον κηδεστὴς ὁ ἐκδιδούς, ἀλλὰ καὶ ὁ λαμβάνων.\\
\francais{Le signe [typographique] : non seulement le Parent est celui qui abandonne [la parole], mais aussi celui qui la [re]prend.}}La scholie ne précise pas à quoi servait ce σημεῖον ; serait-ce une pause ? quoi qu'il en soit, le protagoniste devait certainement faire quelque jeu de théâtre après l'invocation à Zeus du vers 71, identique à celui déjà fait après le premier vers et montrer ainsi aux spectateurs le parallélisme entre les deux quatrains.}
